\documentclass[11pt]{article}
\usepackage[T1]{fontenc}
\usepackage[a4paper,margin=1in]{geometry}
\usepackage{amsmath,amssymb,amsthm}
\usepackage[hidelinks]{hyperref}
\usepackage{microtype}

\newtheorem*{question}{Source open problem}
\newtheorem*{answer}{Later answer}

\setlength{\parindent}{0pt}
\setlength{\parskip}{0.6em}

\title{Literature Status: Diffeological Regularity and an\\
ILB Frobenius Theorem}
\author{Run \texttt{fa\_banach\_001} --- lane 10}
\date{11 August 2026}

\begin{document}
\maketitle

\begin{center}
\fbox{\parbox{0.9\textwidth}{
\textbf{Status: literature already answered.}
The regularity/Frobenius question in arXiv:1710.11076 is explicitly answered
by the same author's later revision arXiv:1607.02636v3.  Cauchy diffeology
and the subset diffeology make the implicit map smooth on a suitable
non-open domain, and Theorem~4.3 proves the corresponding Frobenius theorem.
}}
\end{center}

\section{Original question}

J.-P. Magnot's \emph{On the domain of implicit functions in a projective
limit setting without additional norm estimates}, arXiv:1710.11076
\cite{source}, constructs an implicit map
\[
 u_\infty:D_\infty\longrightarrow F_\infty
\]
on a domain that may fail to be open in the projective-limit topology.  The
domain nevertheless contains the unit ball $\mathcal B$ of an auxiliary
Banach space continuously embedded in the projective limit.

Section~2.2, arXiv PDF page~6, reaches the point where the usual proof of the
Frobenius theorem needs $D_1J$ on $D_\infty$ and asks:

\begin{question}
Classical differentiation applies on open domains, whereas $D_\infty$ need
not be open.  Among the generalized notions of differentiation, which is
best adapted to this setting?
\end{question}

The abstract and introduction identify the same issue as the missing
regularity framework needed for a projective-limit Frobenius theorem.

\section{Separate later answer}

The same author's \emph{On the differential geometry of numerical schemes
and weak solutions of functional equations}, arXiv:1607.02636v3
\cite{answerpaper}, is the separate answering work.  The numbering of the
arXiv identifiers is potentially misleading: v3 was revised on 12 August
2020, after arXiv:1710.11076, and is the version accepted for publication in
\emph{Nonlinearity}.

On PDF page~20 the later paper restates the earlier theorem, says explicitly
that arXiv:1710.11076 left regularity open, and announces that it fills the
gap using Cauchy diffeology.  Its Theorem~4.2 on PDF page~21 proves:

\begin{answer}
There exists a domain $D$ with
\[
 \mathcal B\subset D\subset D_\infty
\]
such that the implicit function $u_\infty$ is smooth when $D$ carries its
subset diffeology.
\end{answer}

The proof realizes $u_\infty$ as the limit of the contraction iterates
$\phi_{i,x}^n(0)$ at every Banach level.  The parameter-to-iterate map is
smooth into the space of Cauchy sequences equipped with the Cauchy
diffeology, and the limit map is smooth by construction.  Passing to the
projective intersection gives the asserted smooth map on $D$.

Theorem~4.3, beginning on PDF page~22 and continuing on page~23, then applies
this result to the path-space implicit equation used in the source.  Under
the levelwise ILB Frobenius compatibility condition, it produces a
diffeological subspace $D$ and a smooth map $J:D\to F_\infty$ satisfying
\[
 D_1J(x,y)=f_i\bigl(x,J(x,y)\bigr)
\]
together with the required normalization on the connected fiber over the
base point.  This is the Frobenius theorem whose classical proof was blocked
in the source.

\section*{Scope}

The later paper selects a precise generalized calculus, proves regularity in
that calculus, and derives the desired Frobenius statement.  It therefore
answers the source's regularity/Frobenius open problem rather than merely
suggesting terminology.

The source has a second, separate closing question about the structure of
$G\mathcal L_\infty/GL(B)$.  The cited theorems do not classify that quotient,
and this packet makes no claim about it.  No new mathematical result is
claimed here; this packet records an author-explicit later-literature answer.

\begin{thebibliography}{9}
\bibitem{source}
J.-P. Magnot,
\emph{On the domain of implicit functions in a projective limit setting
without additional norm estimates},
Demonstr. Math. \textbf{53} (2020), 112--120;
arXiv:1710.11076.

\bibitem{answerpaper}
J.-P. Magnot,
\emph{On the differential geometry of numerical schemes and weak solutions
of functional equations},
Nonlinearity \textbf{33} (2020), 6835--6867;
doi:10.1088/1361-6544/abaa9f; arXiv:1607.02636v3.
\end{thebibliography}

\end{document}
